\section{Cote 136: the weighted PD end retains the tensor obstruction}
\label{sec:136twtor}

\paragraph{Scope.}
This concerns the strict PD functor $\mathsf T$ in
(136.TW2), not all possible common targets or sources.
We prove that its already verified vertical exactness
does not extend to horizontal admissibility: the standard
admissible input $B$ is a counterexample to that particular
functorial construction. The resulting Tor calculation
is the inherited AQ-136-TENSOR calculation, now identified
inside the proposed end. The archive locators remain
cote 136, physical pp.~117--121, footers 116--120,
scans 117, 119 and 121.

\paragraph{A linear end comparison.}
Let $S$ be a commutative ring and let $R=GU$ be the
cofree resolution functor on simplicial $S$-modules.
For simplicial $S$-modules $P,C$, define
\[
 F(P,C)=\int_{[r]\in\Delta}
                  P_r\otimes_S R^{r+1}C.                \tag{136.TO1}
\]
The result is a simplicial module in the original
variable of $C$. Suppose the normalized complex $NP$
is bounded and consists of finitely generated projective
$S$-modules. There is a natural horizontal quasi-isomorphism
\[
                      P\otimes_S C\longrightarrow F(P,C),
                                                               \tag{136.TO2}
\]
where the tensor on the left is the diagonal simplicial
tensor. The same assertion holds for graded $S$ and
graded modules, when ``finitely generated projective''
means a summand of a finite sum of graded free shifts.

\begin{proof}
Every $P_r$ is finitely generated projective by the
Dold--Kan decomposition. Thus its tensor commutes with
the cofree products. We first define (136.TO2) explicitly.
In original simplex degree $n$, a coordinate of
$R^{r+1}C$ is indexed by
\[
 [l_r]\xrightarrow{\alpha_r}\cdots
 \xrightarrow{\alpha_1}[l_0]\xrightarrow{\theta}[n].
\]
Let $\beta:[r]\to[n]$ have the weakly increasing values
\[
 \theta(0),\quad\theta\alpha_1(0),\quad\ldots,
                         \theta\alpha_1\cdots\alpha_r(0).
\]
Send $x\otimes c\in P_n\otimes C_n$ in this coordinate to
\[
           P(\beta)x\ \otimes\
                  C(\theta\alpha_1\cdots\alpha_r)c
                           \in P_r\otimes C_{l_r}.      \tag{136.TO3}
\]
Composition shows that these coordinates obey all end
relations: a coface removes the matching first-vertex
entry, and the last coefficient restriction composes
with the corresponding simplex map; degeneracies repeat
an entry. This is the first-vertex construction of
ARCH-136-TW-CONSTANT with its variable coefficient
carried to the last object of the string. It defines a
natural simplicial $S$-linear map.

We require an exactness observation. Fix the original
simplex degree $n$, and write $A^r=(R^{r+1}C)_n$.
Finite projectivity identifies the end with
\[
 F(P,C)_n=
 \operatorname{Hom}_{c\operatorname{Mod}S}(P^\vee,A^
 \bullet)
 =\operatorname{Hom}_{\operatorname{Ch}^{\geq0}(S)}
                         ((NP)^\vee,N_c A),             \tag{136.TO4}
\]
where $\vee=\operatorname{Hom}_S(-,S)$, $P^\vee$ is
cosimplicial, and $N_c$ is cosimplicial normalization.
The first identity is the finite-projective tensor-dual
identity and the end relation is exactly naturality
of a cosimplicial map. The second is Dold--Kan; the dual
of normalized simplicial chains identifies with normalized
cosimplicial cochains because their degeneracy splittings
are finite and projective.

The augmented complex $C_n[0]\to N_cA$ is split by
the extra codegeneracy after evaluating $n$. Consequently
$N_cA$ is isomorphic, as a cochain complex of modules,
to $C_n[0]$ plus contractible cochain disks. This follows
directly by splitting each differential onto its cycles
with the contracting homotopy. A complex of this form
is injective for degreewise split short exact sequences
of cochain complexes: maps to $M[0]$ are arbitrary maps
from degree zero, and maps to the disk
$M\xrightarrow{1}M$ in degrees $j,j+1$ are arbitrary
maps from degree $j+1$. Such maps extend across a
degreewise split inclusion. Formula (136.TO4) therefore
shows that $F(-,C)$ is exact on degreewise split short
exact sequences of the indicated $P$. This is exactness
in each original simplex, hence in simplicial modules.
The same is true of $-\otimes_S C$.

We now check the projective objects of simplicial modules.
Under Dold--Kan they are summands of sums of $S_0(V)$
and homological disks $D_j(V)$, $j\geq1$, with $V$
projective. We need only finite sums with $V$ finitely
generated. For $NP=S_0(V)$, (136.TO4) is
\[
           \operatorname{Hom}_S(V^\vee,Z^0N_cA)
                       =V\otimes_S C_n.
\]
Here $Z^0N_cA=C_n$ by the augmented resolution.
The map (136.TO3) is the identity under this
identification.

For $NP=D_j(V)$, its dual is the cochain disk in
degrees $j-1,j$. Maps \emph{from} this disk are arbitrary
maps from its degree-$j-1$ module, so (136.TO4) is
\[
                           V\otimes_S(N_cR^{\bullet+1}C)^{j-1}.
                                                               \tag{136.TO5}
\]
The original simplicial module in parentheses is cofree.
Indeed cosimplicial normalization is the intersection
of the kernels of the codegeneracies. Every codegeneracy
$R^{r+1}C\to R^rC$ is the outer right adjoint $G$
applied to a map of degree families: for the first
one this is the formula $G\varepsilon U$ for monad
multiplication, and the others retain an outer $R$.
Since $G$ preserves kernels, their intersection is of
the form $G(M)$. The explicit right-adjoint disk
calculation (136.AT3) shows that $G(M)$ is simplicially
contractible. Tensoring with a finitely generated
projective $V$ commutes with these cofree products and
retains that contraction. Thus (136.TO5) is contractible.
The source $P\otimes C$ is also contractible, since
the simplicial linear contraction of $P$ tensors with
the identity on $C$. This proves (136.TO2) on these
projectives and their finite sums and summands.

Finally a bounded complex of finitely generated
projective modules has a finite, degreewise split
resolution by these projective \emph{chain objects}.
Explicitly, map the sum of $D_j(K_j)$ for $j\geq1$
and $S_0(K_0)$ onto such a homological complex $K$,
using the identity in the top term of each disk and
the differential in its lower term. The map is
degreewise onto and split because every $K_j$ is
projective. Its kernel has finitely generated
projective terms and its largest possible nonzero
degree decreases by one. Iteration terminates.
Apply the two exact functors just verified and use
the natural transformation (136.TO3). The long exact
sequences of original simplicial homology and the
short five lemma inductively extend the projective
case to $P$. This proves (136.TO2).
\end{proof}

\paragraph{ARCH-136-TW-TOR: failure on the standard admissible input.}
Let $k$ be any commutative ring, $S=\Gamma_k(T)$,
and use the standard $B$ over $k$ in (136.TW2).
For every prime $p$ and every integer $a\geq1$,
\[
          H_1\bigl(N\mathsf T(B)_*^{p^a,0}\bigr)
                              \simeq k/pk.               \tag{136.TO6}
\]
Consequently this PD end construction does not send
$B$ to an admissible theory whenever $k$ is nonzero
and is not a $\mathbb Q$-algebra. In particular it
does not rectify the associative comparison over
$\mathbb Z$, even though its vertical cohomology is
correct in every simplex and weight.

\begin{proof}
Write $P=B_*^{*,0}$ and $C=B_*^{*,0}$, as simplicial
graded $S$-modules. The zero-form row of $\mathsf T(B)$
is exactly $F(P,C)$: there is no positive form degree
in either tensor factor of a zero-form element.
Fix a total weight $w$. In simplex degree $r$,
\[
 P_r=S\langle x_1,\ldots,x_r\rangle_{\rm PD}.
\]
Let $P^{\langle w\rangle}_r$ be the free $S$-submodule
on the relative-coordinate PD monomials of degree at
most $w$. It has finite rank over $S$ in every simplex.
All simplicial maps preserve these submodules: they
substitute affine linear expressions in the relative
coordinates with coefficients in $S$, so they cannot
increase their relative degree. Both ends and diagonal
tensors in total weights at most $w$ are unchanged
by replacing $P$ with $P^{\langle w\rangle}$, since
the other coefficient modules have nonnegative weights.

The normalized complex of $P^{\langle w\rangle}$
vanishes above degree $w$. To check this precisely,
filter by relative polynomial degree. Its associated
graded term of degree $b$ is
$S\otimes_k\Gamma_k^b(Q)$, where
$Q=L/k\sum x_i$ and $NQ=k[1]$. In the usual
Dold--Kan coordinates $Q_r=k^r$, degeneracies insert
a zero coordinate. A PD monomial of total exponent
$b$ uses at most $b$ coordinates. If $r>b$, it
omits a coordinate and is therefore a degeneracy
of a monomial in degree $r-1$. Thus the normalized
quotient of $\Gamma^b(Q)$ vanishes in degree $r>b$.
Normalization is exact, so the finite filtration
proves the assertion for $P^{\langle w\rangle}$.
Its normalized terms are finitely generated projective
$S$-modules, since they are summands of its finite
free terms. All statements respect the grading.

The linear end comparison (136.TO2) therefore gives
\[
 H_i\bigl(N\mathsf T(B)^{w,0}\bigr)
       \simeq H_i\bigl(N(B_*^{*,0}\otimes_S B_*^{*,0})^w\bigr)
       \simeq\operatorname{Tor}_i^S(k,k)^w.              \tag{136.TO7}
\]
The last equality and its prime-power evaluation are
the complete inherited AQ-136-TENSOR calculation:
$N(B_*^{*,0})$ is a projective $S$-resolution of $k$;
the diagonal tensor computes its derived tensor; and
$\operatorname{Tor}_1^S(k,k)=I/I^2$ has weight-$p^a$
part $k/pk$, where $I=S_{>0}$. This gives (136.TO6).
For any nonunit prime, these are nonzero in arbitrarily
large weights in form degree zero, violating (136.4).
Every nonzero non-$\mathbb Q$-algebra has such a prime.
\end{proof}

\paragraph{What survives this test.}
The end is still a strict PD algebra and
ARCH-136-TW-VERTICAL remains true. The constant-input
case ARCH-136-TW-CONSTANT remains admissible. The
counterexample concerns their proposed extension to
all variable admissible inputs; it does not affect
the original common-source question, since its
chosen input is the standard admissible theory itself.
